\subsection{Exact normalized solves and their action}

We state precisely the earlier analytic estimates used here.
Lemma~\ref{lin:strip} controls, on unit evolution strips and at a
fixed sufficiently high angular order,
\[
 u,\ u_t,\ u_{\theta\theta},\ E u_{tt},\ \sqrt E\,u_{t\theta}
 \quad\hbox{in }L^2_tH^h_\theta,
 \qquad u,\ u_\theta,\ \sqrt E\,u_t\quad\hbox{as traces}.
\]
The coefficients must have radial and angular diffusion comparable
to \(E\) and \(1\), strictly positive normalized real parts, bounded
normalized derivatives, and mixed coefficient \(o(\sqrt E)\).
Proposition~\ref{lin:global} supplies an invertible augmented block
including the compact observations and multipliers, in uniformly
local growth-weighted versions of these norms.  On a cylindrical
piece the weight has a rate strictly above \(2\), with a positive
buffer above the propagation rate; bounded terminal bands use a
fixed larger rate if necessary.  The conical norms are those of
\Cref{lin:conical}, with the compact commutations used in
\Cref{lin:global}.  The inverse constant is
uniform in \(B\) with these weights retained.  Its physical norm,
and any fixed number of contour or finite-parameter derivatives,
may have fixed polynomial losses.  Exits are pure graph charts
with the outgoing derivative trace prescribed.  The fast Neumann
lift gains \(O(\sqrt E)\) in strip size relative to smooth fixed
angular-order derivative data.  Under conjugation the exit is the
corresponding Robin condition, as in Lemma~\ref{lin:shielding}.

Finally, Lemmas~\ref{gauge:mixed} and \ref{gauge:compact}
give analyticity of the complete prescribed-gauge map,
including the compact transition charts.
On a bounded scaled band, for strip size \(W\) and \(E\asymp B^{-2}\),
its nonlinear remainder and difference are bounded by
\begin{equation}\label{arr:tame-interface}
 C(BW^2+B^2W^3),\qquad
 C(BW+B^2W^2)\|u-\widetilde u\|.
\end{equation}
Second derivatives remain in the strip \(L^2\) factor and first
derivatives in the available traces; these estimates hold at the
full fixed angular order.  Thus they are estimates for the full
operator.  On height bands the corresponding estimate is
Lemma~\ref{gauge:height}.  In the present preliminary construction
only bounded scaled gauge changes are needed.

Here is the exit-data class for that construction.  Let $g_{\rm app}$
be the outgoing derivative of its comparison graph in the prescribed
pure exit variable, after the fixed graph-unit conversion.  If $T$ is
the exit and $\rho_B$ the growth weight of \Cref{lin:global}, permit
exactly the data satisfying
\begin{equation}\label{arr:allowed-traces}
 \rho_B(T)^{-1}
 \|g-g_{\rm app}\|_{\mathcal T^N_{E(T),h}}
       \le c_{\rm exit}B^{-4}.
\end{equation}
For fixed $s$ and contour parameters, use the prescribed real pullbacks
to identify each exit with a fixed boundary circle, and choose the
reference scales $E(T)$ and $\rho_B(T)$ independently of $p$ on the
outer parameter ball.  The varying geometric scales are uniformly
comparable there.  Let $\mathfrak T_B$ be the product of these complex
trace spaces, with norm the maximum of the left-hand norms in
\eqref{arr:allowed-traces}, and write
$\delta g=(g-g_{\rm app})_{\rm exits}$.  The fixed constant
$c_{\rm exit}$ is chosen for the finite collection of homotopies under
consideration, with a fixed factor of room in the contraction
construction.  Along a smooth contour homotopy the same bound holds
for its values, while each required derivative of its normalized
coefficients and data has a fixed polynomial bound.  The natural trace
norm is part of the inverse, so no unproved angular trace is used in
lifting these data.

We also fix the parameter reserve used below.  For constants
$0<c_{\rm in}<c_{\rm out}$ and an exponent $L$, put
\begin{equation}\label{arr:parameter-reserve}
 \mathcal B_s^{\rm out}=\{|p-p_0(s)|<c_{\rm out}|s|^L\},
 \qquad
 \mathcal B_s=\{|p-p_0(s)|<c_{\rm in}|s|^L\}.
\end{equation}
Solves are constructed on the outer ball and all uniform fixed
parameter derivatives are evaluated on the inner ball.  Cauchy's
estimate then costs only a fixed power of $|s|^{-1}$ at each fixed
order.  The finitely many later ball and sector restrictions are chosen
at once; their sectors retain an excess
$\delta_{\rm use}=\vartheta_\delta\delta\asymp D^{-1/2}$ for one fixed
$0<\vartheta_\delta<1$.

\begin{proposition}[Preliminary normalized arrays]\label{arr:preliminary-solves}
Fix \(D\), the bounded scaled portions of the contour families of
Lemma~\ref{arr:contours}, and a fixed lower radius \(r_*>0\).
Let \(p_0(s)\) be a sufficiently high polynomial truncation of
\(\widehat p(s)\).  For some fixed \(L\) and outer ball in
\eqref{arr:parameter-reserve}, for every \(p\in\mathcal B_s^{\rm out}\)
and every choice of outgoing data in \eqref{arr:allowed-traces} at
the artificial exits, the normalized stationary problem on either
contour family has a unique solution in the perturbative neighborhood.
For each fixed contour, the solution operator is jointly holomorphic
in \(p\) and an independent complex trace remainder
\(\delta g\in\mathfrak T_B\) for \(p\in\mathcal B_s^{\rm out}\) and
\(\|\delta g\|_{\mathfrak T_B}<c_{\rm hol}B^{-4}\), with a fixed
\(c_{\rm hol}>c_{\rm exit}\).  Hence a chosen \(p\)-dependent datum gives a holomorphic
solution when \(\delta g(p)\) is holomorphic into that fixed trace
space and remains in the stated trace ball on \(\mathcal B_s^{\rm out}\).
Smooth contour families with smooth prescribed data give
smooth solution families.  Every allowed solution lies in the
growth-weighted augmented ball
\[
 \|(u,\lambda)-(u_{\rm app},\lambda_{\rm app})\|_{X_B}
       \le C_{\rm pre}B^{-4},
\]
where the fixed $C_{\rm pre}$ is chosen before $B$.  In the preceding
strip norms its correction satisfies
\begin{equation}\label{arr:preliminary-bound}
 W(z)\le C_D B^{-4}(1+|z|)^{2.2}.
\end{equation}
Its compact and conical correction is \(O(B^{-4})\).
In particular on \(|z|\asymp B\) its discrepancy from circular data
is \(O(B^{-1.8})\).
The same conclusions hold in one such neighborhood along curtailment,
path, and exit-data homotopies that retain the fixed lower radius and
\eqref{arr:allowed-traces}, with the just stated regularity of the data
for the corresponding family-regularity conclusion.  This includes the first fixed-radius
height turn, but not the subsequent shrinking height turns.  Fixed
finite parameter derivatives on $\mathcal B_s$ and the prescribed
smooth contour derivatives have only polynomial losses in \(B\) and
\(|s|^{-1}\).
\end{proposition}

\begin{proof}
Use the exact family \(S_b\) as base, with its multiplier and the
first multiplier jet in the comparison pair.  Remove the first-order
transverse tilt by analytic rotations of the individual ends,
cut off on fixed compact collars.  On the real incoming part add
the linear opening jet and the nonlinear part of \(R_c\), starting
the latter beyond a fixed end entrance.  Before changing the
spatial contour, cut off the small noncircular remainders on a
bounded logarithmic band and use \(R_c\) alone on the complex part.
All cutoffs are prescribed functions on the real parameter domain.
They are never evaluated at an unknown complex spatial argument.

Here is the size of the residual.  On the compact and conical
pieces the constant and linear terms have been solved, so it is
\(O(B^{-4})\), including the observation error.  On a cylindrical
piece the no-axial-derivative circular equation vanishes exactly
on \(R_c\).  The base error is \(O(z^{-2+\eta})\); the nonleading
part of the first opening jet is \(O(B^{-2}z^\eta)\).
Their interactions beyond first order have size at most
\(CB^{-4}(1+|z|)^{2+\eta'}\).
The axial curvature terms carry \(z^{-2}\); expansion where
\(|\beta z^2|\) is small gives the same bound.
Choose the fixed losses so that \(\eta'<0.06\).  On the final
bounded logarithmic bands, direct differentiation of \(R_c\)
and of the prescribed cutoff gives \(O(B^{-2+\eta'})\).
Thus the output residual is bounded by
\[
 C_D B^{-4}(1+|z|)^{2.06}.
\]
Small compact corrections enforce the observations exactly.
The exit derivative is that of the comparison profile, so its
initial derivative residual is zero.

Apply Proposition~\ref{lin:global} with rate \(2.2\), keeping the strict
buffer above \(2.06\), and with the bounded-band weights specified
above.  Its hypotheses hold: before the bounded outer part the
coefficients are small perturbations of the cylindrical family;
on that part Lemma~\ref{arr:contours} gives the diffusion signs,
and the reference remains transverse with nonsingular drift.

We verify the index hypothesis separately, using a homotopy of
operators on fully stopped domains. Straightening an adverse
geometric bypass would cross the pole of the circular profile;
the following argument does not do so.

Fix $B$ and finite conical stopping sections, and pull every chart
back to its prescribed real parameter domain. The scalar graph
variations and scalar normal sources describe the same complex
normal line of the reference immersion. Identifying a graph
variable with its normal component multiplies the unknown and
the source by the same nonzero graph-speed factor. This is
conjugation of the operator and changes only lower-order terms.
Its principal tensor is therefore a single geometric tensor $A$.
The local drift divisions give smooth nonzero row multipliers
$m_i$ for which $\operatorname{Re}(m_iA)$ is positive definite,
with the sign chosen for minus the second-order part. This is
the strict radial and angular accretivity and small mixed
coefficient on the cylindrical and height charts. It follows
from closeness to the real reference on the core and its collars,
and from Lemma~\ref{lin:conical} on conical pieces. The spatial
coordinate changes on the pulled-back domains are real, so
transform the tensors by real congruence and preserve positivity.

Choose a nonnegative real partition of unity $\chi_i$ subordinate
to these charts and define $m=\sum_i\chi_i m_i$. Then, for every
nonzero real covector $\xi$,
\[
 \operatorname{Re}\bigl(mA(\xi,\xi)\bigr)
 =\sum_i\chi_i\operatorname{Re}\bigl(m_iA(\xi,\xi)\bigr)>0.
\]
Consequently $m$ is nonzero everywhere. Multiplication of the
entire source row, including the compact multiplier columns, by
$m$ is a range isomorphism and produces one globally accretive
principal tensor. This avoids imposing an unsupported common
phase on the different local normalizing factors.

The graph-unit identification may write an actual exit condition
as $(\partial_t+q)u=0$. Choose a smooth function $F$ supported
in a collar of that exit, with $F_t=-q$ on its boundary, and put
$u=e^Fv$, multiplying the source by $e^{-F}$ as well. The
boundary condition becomes exactly $v_t=0$, while the principal
tensor is unchanged. Disjoint collars allow this at every exit;
entrance Dirichlet data stay zero. Fix all these isomorphisms
before the homotopy. For the compact index calculation take the
ordinary $H^2$ domain with these homogeneous boundary conditions
and range $L^2$, together with the finite multiplier and observation
spaces. These spaces remain fixed along the homotopy. Positive smooth
weights at a fixed stop give equivalent norms at this order; they
do not identify the higher angularly commuted spaces with an
isotropic higher Sobolev scale.

Choose a real Riemannian metric which is a product in the exit
collars, so that $\partial_t$ there is a nonzero multiple of its
normal derivative. Convexly interpolate the globally normalized
principal tensor to its positive inverse metric, and interpolate
the lower coefficients to those of $-\Delta+1$. Every intermediate
principal tensor has positive real part. The scalar boundary
symbol has one decaying normal root for each nonzero real
tangential covector. This root cannot be zero, since the
tangential quadratic form has positive real part. Thus the fixed
transverse derivative condition stays complementing; so does
Dirichlet data. The endpoint $-\Delta+1$ with these mixed
Dirichlet/Neumann conditions is invertible by its coercive variational
form and boundary regularity, and has index zero.
Keep the transformed compact observations and multiplier
columns fixed along this homotopy. Their equal-dimensional
augmentation has index zero too.

This proves the required index assertion on the actual fully
stopped domains. Boundary regularity makes its kernel smooth, so
the weighted estimate applies at the actual geometric operator and
gives bijectivity on the compact $H^2$--$L^2$ realization.
For smooth data the solution obeys the full angularly commuted
target estimates; smooth approximation then gives the inverse
from $Y_B$ to $X_B$, with the same bound. Exhaustion of the
conical ends then proceeds exactly as in Proposition~\ref{lin:global};
the artificial compact homotopy is not used on the infinite ends.

The resulting linear correction satisfies
\eqref{arr:preliminary-bound}.  On every outer strip its
nonlinear Lipschitz factor is at most
\[
 C_D(BW+B^2W^2)=O_D(B^{-0.8}),
\]
and on earlier strips the local axial scale is at most \(C_DB\),
so the factor is no larger.  Fixed compact and conical pieces
have the ordinary small perturbation factors.  For \(B\) large
the correction map therefore preserves a fixed multiple of
\eqref{arr:preliminary-bound} and contracts there.
This proves existence and uniqueness in that neighborhood.

This application needs the uniform weighted inverse.  A merely
polynomial bound for the original correction would not imply
the small factor just displayed.  Polynomial losses are used
only after the exact perturbative solution has been constructed.

The construction is uniform over the compact collections of normalized
contours for fixed \(D\) and fixed lower radius.  The outgoing lift of
\eqref{arr:allowed-traces} has $X_B$ size at most $Cc_{\rm exit}B^{-4}$.
Choose a fixed $c_{\rm hol}>c_{\rm exit}$ and $C_{\rm pre}$ to include
that larger trace ball; the same
contraction then works on an open neighborhood of the stated closed
data class for large $B$.  The full stops use
$g=g_{\rm app}$.  Curtailment and smoothing use the comparison derivative
of the same local reference at their moving exit; at the common pre-pole
endpoint both problems have the same path, scalar exit unit, and
comparison derivative.  The real stopped family below checks the
consecutive-stop data separately.  At the initial fixed-radius height
transition, \Cref{jump:compatibility} verifies the natural trace size
$O(B^{-2.8})$, below the available outer strip size
$B^{-4}\rho_B(T)\asymp B^{-1.8}$, and places the solve in the same
preliminary ball.  Its separate refined inverse covers the later
shrinking paths and their trace interpolations.
For each fixed contour, choose a bounded complex-linear outgoing trace
lift on the identified spaces, supported outside the observation core.
After subtracting this lift, the unknown has homogeneous exits, and its
equation is jointly analytic in the unknown, \(p\), and the independent
remainder \(\delta g\).  Uniform contraction on the larger trace ball
gives the asserted joint Banach-space holomorphy.  The actual full-stop
and curtailment data are comparison derivatives of the analytic
reference profiles at prescribed real path parameters; their remainders,
and affine interpolations with coefficients independent of \(p\), are
holomorphic in \(p\).  The consecutive real stops have the same
property.  Proposition~\ref{jump:compatibility} checks it for the
fixed-radius comparison interpolation and for the later actual
shorter traces in the refined solver.  Cauchy estimates on
$\mathcal B_s$ cost \(C_D|s|^{-Lj}\) at each fixed derivative order \(j\).
Differentiating a smooth contour homotopy uses the same
linearized inverse and the stipulated differentiated coefficient and
data bounds.
This proves the final derivative assertion.

There is no requirement that \(b=O(B^{-2})\): typically
\(b=O(s)=O(B^{-2/k})\).  It belongs to a fixed small analytic
neighborhood of \(S_b\), which is an exact base family.
Only the openings have size \(O(B^{-2})\).
\end{proof}

Let \(G_\pm(s,p)\) be the Gaussian action integrated all the way
to the indicated artificial exits of these solves.  The \(s\)
argument records the contour and exit prescription; it is
distinct from the holomorphic finite parameter \(p\).

\begin{lemma}[Endpoint first variation]\label{arr:first-variation}
At fixed compact observation values, a differentiable family of
the preceding exact normalized solves satisfies
\[
 \frac{dG}{d\tau}
   =\sum_{\text{exits}}\text{boundary position and momentum terms}.
\]
This algebraic identity holds for prescribed artificial Neumann data
satisfying \eqref{arr:allowed-traces}, and for the refined class later
proved in \Cref{jump:compatibility}.  No natural boundary condition
or vanishing boundary momentum is required.
For the preliminary fixed-radius, bounded-scaled families, if the moving exits satisfy
\(\operatorname{Re}z^2\ge H\), and the normalized shapes, exit
data and their derivatives have polynomial bounds in \(B\),
then, after any fixed finite parameter differentiations,
\begin{equation}\label{arr:endpoint-estimate}
 \left|\frac{dG}{d\tau}\right|
 \le C_D B^M e^{-H/4+o(B^2)}.
\end{equation}
Consequently, on slightly smaller parameter balls and sectors,
\begin{equation}\label{arr:dbar}
 |\bar\partial_s\partial_p^\alpha G_\pm(s,p)|
 \le C_{\alpha,D}|s|^{-M_\alpha}e^{-D/(5|a|)}
\end{equation}
for each fixed \(\alpha\).
\end{lemma}

\begin{proof}
Write the pulled-back area density as
\(\mathcal L(X,DX)=e^{-X\cdot X/4}\sqrt{\det g}\).
For a domain with moving boundary, differentiation and one
integration by parts give
\[
 \delta G
  =\int_\Omega \mathcal Q(X)\cdot\delta X
    +\int_{\partial\Omega}\Pi_\nu\cdot\delta X
    +\int_{\partial\Omega}\mathcal L\,v_{\partial\Omega},
 \qquad \Pi_\nu=\mathcal L_{DX}\nu .
\]
The equation and the compact constraints turn the interior term
into \(\lambda\cdot\delta\mathsf P(p)\).
Tangential changes contribute only to the boundary: algebraic
reparametrization covariance gives
\(\mathcal Q(X)\cdot DX=0\).
These identities continue bilinearly to the prescribed complex
charts.  The observation charts are fixed on their compact
support, so no moving-coordinate term is hidden in
\(\delta\mathsf P\).  At fixed \(p\) the interior term is zero.
Prescribed Neumann data determine a solve; they do not set
\(\Pi_\nu\) equal to zero, and the displayed boundary terms
are retained.

For those preliminary families, the area and momentum prefactors, endpoint velocities and
solution variations have polynomial bounds by
Proposition~\ref{arr:preliminary-solves}.
Along these fixed scaled contours,
\(\operatorname{Re}(X\cdot X)=\operatorname{Re}z^2+o(B^2)\);
bounded transverse radii and the small relative graph
corrections account for the error.  This proves
\eqref{arr:endpoint-estimate}.
At the full stop \(H=DB^2\), absorb the polynomial factors
and the \(o(B^2)\) loss into the strict margin from \(1/4\)
to \(1/5\).  Variation of \(s\) at fixed \(p\) changes only
the domain, contour and exit data.  Applying the identity
to its two real derivatives gives \eqref{arr:dbar}.
Finite \(p\)-derivatives follow by Cauchy estimates on the
slightly larger ball.  The same reasoning treats changes of
exit prescriptions at fixed observations.
The refined family uses the same algebraic identity, with its separate
shape and Gaussian-height estimates in \Cref{jump:endpoints}; its trace
class alone is not used to infer \eqref{arr:endpoint-estimate}.
\end{proof}

In particular, on a real evaluation ray the adverse leading poles
\(Z_{j,\mathrm{lead}}^2=(-c_ja)^{-1}\) are positive, and the scale
\eqref{arr:A0} is their minimum divided by four.
Curtail the two contours at
\(\operatorname{Re}z^2=\tfrac34 Z_{j,\mathrm{lead}}^2\) and use
\eqref{arr:affine-homotopy} to make their boundary problems
identical.  Their solutions then agree by perturbative uniqueness.
Integration of \eqref{arr:endpoint-estimate}, with fixed
polynomial losses, proves
\begin{equation}\label{arr:fixed-parameter-comparison}
 |\partial_p^\alpha(G_+-G_-)(s,p)|
 \le C_{\alpha,D}|s|^{-M_\alpha}
          e^{-(0.7-o(1))A_{\mathrm{lead}}}.
\end{equation}
Here and below the notation permits replacing \(0.7-o(1)\)
by any slightly smaller fixed number for all sufficiently
small \(s\).  The actual endpoint exponent tends to \(0.75\).
Nonadverse and smaller-order openings may retain common
farther exits.  For parameters real to leading order only,
the same argument has \(1+o(1)\) errors in the leading
Gaussian heights and remains valid.  It does not require
that the two arrays be conjugates.

\subsection{Full-disk approximants and Gevrey remainders}

The exact contour solves already have small endpoint errors.  We next
relate their actions to the formal action, with estimates at every Taylor
order.  Real contours stopped at increasing radii give holomorphic
approximants on shrinking parameter disks.  Their exponentially small
successive differences are what produce the factorial bounds below.
The relation between Gevrey remainders and exponential flatness is
classical \cite[Sections~1.3--1.6]{MR1992}; the proof here constructs
the required approximants for these normalized stationary problems.

\begin{lemma}[Real stopped polydiscs]\label{arr:real-stops}
There are a fixed small complex $b$-neighborhood, $d_0>0$, and $H_0$
with the following property.  For every independent stop length
$H\ge H_0$, stop each cylindrical end on its real path at height
comparable to $H$, leaving the conical ends infinite.  For
\[
 b\text{ in that neighborhood},\qquad |x|<d_0H^{-2},
\]
the normalized comparison problem has a solution $U_H(b,x)$,
holomorphic in these finite parameters.  In the analogue of $X_B$
with $B$ replaced by $H$, its correction from the base-plus-first-jet
comparison pair is at most $C H^{-4}$.  Thus the cylindrical strip
bound is $C H^{-4}(1+z)^{2.2}$ and the compact and conical correction
is $O(H^{-4})$.  The paths and cutoffs are real and independent of the
finite parameters, and the outgoing comparison derivatives are
holomorphic in them.  For real parameters the stopped surface is real.

Let $F_H(b,x)$ be its Gaussian action.  It is holomorphic and uniformly bounded on a
fixed smaller polydisc.  If $H\le H'\le\Lambda H$ for fixed $\Lambda$,
then on the common smaller polydisc
\begin{equation}\label{arr:real-stop-difference}
 |F_{H'}-F_H|\le C_\Lambda e^{-c_\Lambda H^2}.
\end{equation}
For every fixed $0<\nu<1$, the real $U_H$ is a small normal graph over
$S_b$ on $|Y|\le2H^\nu$ for large $H$.  Its action outside the part
over $|Y|\le H^\nu$ is at most
$C H^d e^{-cH^{2\nu}}$.
\end{lemma}

\begin{proof}
The length $H$ is a parameter of this boundary problem, not
$|s|^{-k/2}$.  Since $\beta(b,x)$ is linear in $x$ with bounded
coefficients, choosing $d_0$ small keeps $|\beta|H^2$ uniformly small.
For real parameters the chosen circular branch has positive radius.
On the complex polydisc it remains in a fixed small neighborhood of
that positive real comparison, so its denominator units and the strict
normalized diffusion bounds on the real path persist throughout the stop.  Use the
same base, first opening jet, and prescribed outer cutoff as in
\Cref{arr:preliminary-solves}.  On the outer cutoff band the base
noncircular remainder is $O(H^{-2+\eta})$, and the nonleading first
jet is $O(H^{-2+\eta})$.  With $\eta<0.06$, their cutoff residual is
bounded by $CH^{-4}(1+z)^{2.06}$ there.  Before that band the solved
constant and linear terms give the same bound; on the compact and
conical pieces it is $O(H^{-4})$.  The weighted inverse and contraction
from the preliminary proof therefore apply with $H$ in place of $B$.
This argument includes $x=0$.  In that case the stopped base can differ
from $S_b$ because the noncircular tail was cut off before the exit.

The data and the operator are analytic in $(b,x)$ on the displayed
polydisc, and uniform contraction gives the asserted holomorphy.
For consecutive stops move the real endpoint and always use the
comparison derivative in its current pure graph coordinate.  Any
interpolation needed to identify the two outer cutoffs changes a
smooth derivative by $O(H^{-2+\eta})$.  Its natural Neumann norm is
$O(H^{-3+\eta})$, since $\sqrt E\asymp H^{-1}$ there.  This is below
the permitted trace size $H^{-4}\rho_H(T)\asymp H^{-1.8}$, so the
whole homotopy consists of exact normalized solves.  Along it
$\Re(X\cdot X)\ge cH^2$ at the exit, and all position, momentum,
and variation factors grow at most polynomially.  Endpoint first
variation proves \eqref{arr:real-stop-difference}; changing $c$ absorbs
the polynomial factors and the fixed ratio $\Lambda$.  On the conical
ends the small complex link graphs retain
$\Re(X\cdot X)\ge c|Y|^2-C$, and their first traces give polynomial
area bounds.  The resulting integrable Gaussian dominates the holomorphic
density uniformly on a smaller parameter polydisc.  Differentiation
under that integral gives holomorphy of $F_H$ and its uniform action bound.

For the final assertion take real parameters.  On $z\le3H^\nu$ the
opening displacement is $O(H^{-2}z^2)$, an end tilt contributes
$O(H^{-2}z)$, and the correction is $O(H^{-4}(1+z)^{2.2})$ in the
strip and first-trace norms.  These quantities and their physical first
derivatives tend to zero for $\nu<1$.  On a conical end the first opening
jet with compact exterior source belongs to the conical entrance domain
by the inverse identification in the family proof.  Its link displacement
is $O(H^{-2})$, and the correction has link norm $O(H^{-4})$.
The conical first-derivative traces therefore give the same smallness
on $|Y|\le3H^\nu$.  This larger region remains inside every cylindrical
stop for large $H$, since $\nu<1$; the conical ends are infinite.
Choose the real projection branches there.  Their inverse projections
over $|Y|\le2H^\nu$ stay in the larger region, because the graph
displacements and slopes tend to zero while the radial margin is $H^\nu$.
The uniform real tubular neighborhood of $S_b$ then puts the graph
descriptions in a common normal graph on $|Y|\le2H^\nu$.
On the remaining real cylindrical
pieces $|X|\asymp z$ and the area density has polynomial bounds from
the first traces; on the conical pieces $|X|$ is comparable to their
radial coordinate and their area density is polynomial.  Integrating
the Gaussian on those tails proves the claimed action bound.
\end{proof}

\begin{lemma}[Geometric telescoping]\label{arr:gevrey}
For any sufficiently high real-coefficient polynomial path \(p_0(s)\) with the
specified opening jet and any sufficiently large fixed \(\ell\),
the arrays \(G_\pm(s,p_0(s)+s^\ell h)\), uniformly for \(h\)
in a fixed complex ball, have the formal stationary expansion
of Gevrey order \(1/k\).  More precisely, if that expansion is
\(\sum_{n\ge0}f_n(h)s^n\), then
\begin{align}
 \|f_n\|&\le C A^n\Gamma(1+n/k),\label{arr:gevrey-coeff}\\
 \left\|G_\pm(s,p_0(s)+s^\ell h)
      -\sum_{n<N}f_n(h)s^n\right\|
   &\le C A^N\Gamma(1+N/k)|s|^N
       \quad(N\ge0).\label{arr:gevrey-remainder}
\end{align}
The norm is the supremum on a fixed smaller \(h\)-ball.
These statements hold also after any fixed finite number
of parameter derivatives.  The constants can depend on \(D\).
\end{lemma}

\begin{proof}
Take the real stopped problems of Lemma~\ref{arr:real-stops} with
\(R_j=R_0\Lambda^j\), \(\Lambda>1\), and set
\[
 F_j(s,h)=F_{R_j}(p_0(s)+s^\ell h).
\]
There are full complex \(s\)-disks
\[
 |s|\le r_j=dR_j^{-2/k}
\]
on which this construction is holomorphic in \((s,h)\), uniformly
for $h$ in a fixed complex ball.  Indeed $\ell$ is chosen at least $k$,
so \(x=O(s^k)=O(R_j^{-2})\), and $b$ remains in the fixed neighborhood.
Choose $d$ small enough for the stopped polydisc.  The independent
length formulation includes the center $s=0$; it does not substitute an
infinite value for $B$ there.

For consecutive radii the common disk is the smaller one.
Equation~\eqref{arr:real-stop-difference} gives
\[
 \|F_j-F_{j-1}\|_{|s|\le r_j}\le C e^{-cR_j^2};
\]
the fixed ratio of successive radii has merely changed \(c\).
Polynomial factors are absorbed by decreasing \(c\).
Set \(D_j=F_j-F_{j-1}\) for \(j\ge1\), and \(D_0=F_0\).
Cauchy's estimate yields
\[
 \|[s^n]D_j\|
 \le C e^{-cR_j^2}d^{-n}R_j^{2n/k}.
\]
For \(t_j=R_j^2\), reserve \(e^{-ct_j/2}\) for the convergent
geometric-scale sum, and maximize the remaining expression:
\[
 \sup_{t>0}t^{n/k}e^{-ct/2}
    =\left(\frac{2n}{cke}\right)^{n/k}.
\]
It follows that
\begin{equation}\label{arr:moment-sum}
 \sum_{j\ge1}e^{-cR_j^2}R_j^{2n/k}
       \le C A^n\Gamma(1+n/k).
\end{equation}
Thus
\(\widetilde f_n=[s^n]F_0+\sum_{j\ge1}[s^n]D_j\)
exists as a holomorphic function of $h$ and satisfies the coefficient
bound in \eqref{arr:gevrey-coeff}.  At this point it is the coefficient
series of the stopped telescope; its identification is proved below.

For the all-orders remainder, fix \(0<\vartheta<1\) and take
\(J=J(s)\) maximal with \(|s|\le\vartheta r_J\).
Then \(R_J^2\asymp |s|^{-k}\).
Writing \(T_{N-1}\) for the Taylor polynomial through degree
\(N-1\), the exact identity is
\begin{equation}\label{arr:telescoping-remainder}
 F_J-\sum_{n<N}\widetilde f_ns^n
 =\sum_{j=0}^J(D_j-T_{N-1}D_j)
       -\sum_{j>J}T_{N-1}D_j .
\end{equation}
For \(j\le J\), the Taylor remainder is at most
\[
 \frac{C e^{-cR_j^2}}{1-\vartheta}
       |s|^N r_j^{-N},
\]
with the ordinary analytic bound for \(D_0\).
For \(j>J\), put \(x_j=|s|/r_j>\vartheta\).
For every \(N\ge1\),
\[
 \sum_{n<N}x_j^n\le N\vartheta^{-N}x_j^N.
\]
Use \eqref{arr:moment-sum} on both sums in
\eqref{arr:telescoping-remainder}, absorbing
\(N\vartheta^{-N}\) into \(A^N\).  This proves the stated
remainder for every \(N\), including orders beyond optimal
truncation.

We now identify the coefficient series using only real graph
comparisons.  Fix an order $n$ and choose $M$ with $k(M+1)>n$.
Choose $0<\nu_M<1$ so small that the degree-$M$ polynomial height and
its first two derivatives are $o(1)$ on $|Y|\le2H^{\nu_M}$ whenever
$|x|\le CH^{-2}$.  This is possible by its fixed polynomial growth;
for example, if that growth has degree $D_M$, take
$\nu_MD_M<1/2$.  On real parameters Lemma~\ref{arr:real-stops} puts
$U_H$ in the same small normal graph there.  Its normalized residual
and observation error vanish: the real change of graph variable
transforms the scalar source row with the graph speed, while the compact
observations and the multiplier $\lambda_H$ retain their fixed
variational normalization.  Fix the real Gaussian exponent $a_G=1$.
The small graph neighborhood can be chosen with retained density exponent
$c_G>1/2$, so $a_G<2c_G$.  In the reference comparison lemmas use their
weight $a=a_G$.  Applying
\Cref{ref:real-comparison,ref:finite-taylor} with outer radius
$R=2H^{\nu_M}$ and $\vartheta=1/2$ gives
\[
 \|\chi_R(u_H-v^{(M)})\|_{\mathcal H^2_{a_G}}
       +|\lambda_H-\lambda_M|
 \le C_M|x|^{M+1}+H^{d_M}e^{-c_MH^{2\nu_M}}.
\]
Here the annular commutator is controlled by the real first-trace
bounds and the $C^2$-small Taylor graph.  The real action-tail bound
of Lemma~\ref{arr:real-stops} and \eqref{ref:real-action-estimate}
then yield
\begin{equation}\label{arr:stopped-taylor-value}
 |F_H(p)-f_M(p)|
 \le C_M|x|^{M+1}+H^{d_M}e^{-c_MH^{2\nu_M}}.
\end{equation}
The Taylor graph in this comparison is used only on the inner domain.
The full $f_M$ is the integral of its finite polynomial coefficient
densities, whose tails were included in the estimate.  Thus the argument
neither requires a global uncut polynomial immersion nor changes real
normal graphs by a complex spatial reparametrization.

Apply \eqref{arr:stopped-taylor-value} with $H=R_{J(s)}$ on $s>0$ and
$h$ in a fixed smaller real ball.  The balance $H^2\asymp s^{-k}$
and $|x|=O(s^k)$ make the right side
$O(s^{k(M+1)})$ plus a function smaller than every power of $s$,
uniformly on this real ball.  For the fixed $M$, the function
$f_M(p_0(s)+s^\ell h)$ is an ordinary analytic function of $(s,h)$
on a sufficiently small neighborhood.  The formal composition of
$\mathcal F$ has the same coefficients through order $n$: every
omitted $x$ monomial has formal $s$-valuation at least $k(M+1)$.
This valuation statement does not assign a numerical bound to the
omitted formal tail.

The already proved telescoping remainder, with $N=n+1$, is uniform
on the smaller complex $h$-ball, even though $J(s)$ jumps.  Compare it
with the ordinary Taylor expansion of $f_M(p_0(s)+s^\ell h)$ and
\eqref{arr:stopped-taylor-value}.  Inductively subtract the lower
identified powers, divide by $s^n$, and let $s\downarrow0$.
This identifies $\widetilde f_n(h)$ with the coefficient $f_n(h)$ of
the formal stationary composition on the real ball.  Both sides are
holomorphic in $h$, so the identity extends to the complex ball.
Since $n$ was arbitrary, the stopped coefficient series is exactly
the formal stationary action.  The weaker cutoff exponential in
\eqref{arr:stopped-taylor-value} is used only for this identification;
the $e^{-cR_j^2}$ endpoint telescope supplies all Gevrey bounds.

The sector array and \(F_J\) can be curtailed to a common
inner real-path region at radius \(c_DB\).
Their boundary problems are connected by the homotopies
already constructed, at fixed observations.
Endpoint first variation gives
\[
 \|G_\pm-F_J\|\le C_D e^{-b_D/|s|^k},\qquad b_D>0.
\]
For every \(N\),
\[
 e^{-b_D/|s|^k}
 \le \left(\frac{N}{b_Dke}\right)^{N/k}|s|^N
 \le C A_D^N\Gamma(1+N/k)|s|^N .
\]
This transfers the all-orders remainder.
Finally, Cauchy's estimate on nested fixed \(h\)-balls gives these
bounds for derivatives in the scaled variable \(h\).
Derivatives in the original parameter \(p\) additionally require
fixed division; their proof is given after Lemma~\ref{arr:division},
in \eqref{arr:unscaled-remainder}.
The gamma weight is equivalent, up to exponential factors in \(N\),
to \((N!)^{1/k}\).
\end{proof}

\begin{lemma}[Fixed divisions and regular implicit operations]
\label{arr:division}
The class of arrays in Lemma~\ref{arr:gevrey}, with their
full-disk exponentially telescoping approximants, is stable
under finite sums, products, analytic operations on a fixed
bounded neighborhood, and regular finite-dimensional implicit
solution.  It is also stable under division by a fixed power
\(s^q\) when its first \(q\) formal coefficients vanish
identically on the parameter ball.  Division means the
quotient of actual sector values.
A formally zero result satisfies \(Ce^{-c/|s|^k}\) for
some \(c>0\).
\end{lemma}

\begin{proof}
Sums, products and a fixed analytic operation preserve
exponential telescoping by their bounded derivatives on
slightly smaller fixed neighborhoods.
For division the exact remainder identity is
\[
 R_N(s^{-q}F)=s^{-q}R_{N+q}(F).
\]
The ratio of gamma weights with the fixed shift \(q\)
has at most polynomial growth in \(N\), which is absorbed
in a larger exponential base.

The disk approximants need a correction:
formal vanishing of the limiting coefficients does not
make \(F_j/s^q\) holomorphic.  Use instead
\begin{equation}\label{arr:low-jet-subtraction}
 \widetilde F_j=\frac{F_j-T_{q-1}F_j}{s^q}.
\end{equation}
On a fixed smaller disk their differences are bounded by
\[
 C r_j^{-q}e^{-cR_j^2}\le C e^{-c'R_j^2}.
\]
Each removed low coefficient has limit zero and is therefore
an exponentially small tail of the coefficient telescoping
sum.  At \(J(s)\), its value after division by \(s^q\)
has only a fixed polynomial loss, so remains exponentially
small.  Thus \eqref{arr:low-jet-subtraction} approximates the
actual quotient with the required accuracy.
Uniform formal vanishing on the whole parameter ball is
essential in this argument.

For implicit solution suppose the normalized equation is
\(H(s,z)=0\), with leading equation \(H_0(z)=0\), root \(z_0\),
and invertible \(A=D_zH_0(z_0)\).
Take the corrected disk approximants \(H_j\).
Their centered telescoping gives, for any fixed
\(0<\vartheta<1\),
\[
 \|H_j(s,\cdot)-H_0\|_{|s|\le\vartheta r_j}
       \le C r_j+C e^{-c'R_j^2}.
\]
In fact the constant-coefficient tail is exponentially small,
while the nonconstant parts are bounded by \(|s|\) times
the convergent sum
\(\sum_{\nu\ge1}e^{-cR_\nu^2}r_\nu^{-1}\), together with
the fixed initial analytic term.  Cauchy's estimate on a
smaller \(z\)-ball gives the same convergence for the
first \(z\)-derivative.  Shrink that fixed ball and the
disks so that
\[
 \Psi_{j,s}(z)=z-A^{-1}H_j(s,z)
\]
is a contraction with a common constant \(q_*<1\) and maps
that ball into itself.  Its roots \(z_j(s)\) are holomorphic,
and subtraction of their equations gives
\[
 \|z_j-z_{j-1}\|
 \le\frac{\|A^{-1}\|}{1-q_*}\,
          \|H_j-H_{j-1}\|.
\]
They therefore telescope exponentially on full disks.
Lemma~\ref{arr:gevrey} supplies their coefficient and
all-orders remainder bounds.  The same contraction estimate
compares the actual sector root with the stopped approximant.
This proves implicit closure with remainders.
If every formal coefficient is zero, choose
\(N=\lfloor c_0|s|^{-k}\rfloor\) in the remainder bound,
where \(c_0>0\) is small.  Stirling's bound gives
\(Ce^{-c/|s|^k}\).
\end{proof}

\paragraph{Unscaled parameter derivatives.}
To complete the derivative assertion of Lemma~\ref{arr:gevrey}, fix
a multi-index \(\alpha\), put \(q=\ell|\alpha|\), and write
\[
 F(s,h)=G(s,p_0(s)+s^\ell h),\qquad
 H=\partial_h^\alpha F
   =s^q(\partial_p^\alpha G)(s,p_0(s)+s^\ell h).
\]
Only the undifferentiated Gevrey statement is used initially.
Cauchy's estimate on nested fixed \(h\)-balls gives the all-orders
bounds for \(H\), and exponential telescoping and sector comparison
for \(H_j=\partial_h^\alpha F_j\).
The ordinary formal jet \(\mathcal F\in\mathbb R\{b\}[[x]]\)
has no negative powers after substitution: every occurrence of
\(h\) carries \(s^\ell\).
Thus, if \(\widehat F=\sum f_n(h)s^n\), then
\(\partial_h^\alpha f_n\equiv0\) for \(n<q\), identically on the
whole parameter ball.  This includes mixed base and opening
derivatives, and follows from the formal chain rule, independently
of any unscaled derivative bound.

Apply the fixed-division part of Lemma~\ref{arr:division}.
The holomorphic disk approximants for the actual quotient \(s^{-q}H\)
are
\[
 K_j=\frac{H_j-T_{q-1}H_j}{s^q},\qquad
 \|K_j-K_{j-1}\|_{|s|\le\vartheta r_j}
       \le C_\alpha r_j^{-q}e^{-cR_j^2}
       \le C'_\alpha e^{-c'R_j^2},\quad 0<c'<c .
\]
For \(q=0\) no subtraction is made.
To check the actual sector values, uniform formal vanishing gives
\([s^m]H_J=-\sum_{j>J}[s^m](H_j-H_{j-1})\) for \(m<q\).
At the balanced index
\(\vartheta r_{J+1}<|s|\le\vartheta r_J\), these finitely many
tails and \(H-H_J\) have a common positive exponential decay
constant \(b\), and therefore give
\[
 \|s^{-q}H-K_J\|
 \le |s|^{-q}\|H-H_J\|
       +\sum_{m<q}|s|^{m-q}\|[s^m]H_J\|
 \le C_\alpha e^{-b'/|s|^k},\qquad 0<b'<b .
\]
The decreased exponent absorbs only fixed powers of \(|s|^{-1}\);
individual stopped approximants were not assumed divisible.
Finally, with \(R_N\) denoting the remainder after degree \(N-1\),
\begin{equation}\label{arr:unscaled-remainder}
 \|R_N(s^{-q}H)\|
 =|s|^{-q}\|R_{N+q}(H)\|
 \le C_\alpha A_\alpha^N\Gamma(1+N/k)|s|^N
 \quad(N\ge0).
\end{equation}
Indeed
\(\Gamma(1+(N+q)/k)/\Gamma(1+N/k)\le C_{q,k}(N+1)^{q/k}\);
a larger fixed exponential base absorbs this ratio.
The coefficient bound follows by the same fixed shift.
This proves the assertion for every order, including beyond optimal
truncation, using a division proof that requires no unscaled
derivative closure.


\subsection{Replacing a divergent critical arc}

The formal critical arc need not converge, so it cannot simply be
substituted into an actual contour action.  We select finitely many
equations with maximal rank, approximate their free coordinates by
polynomials, and solve for the remaining coordinates.  The next lemma
shows algebraically why this replacement retains every original formal
constraint; the proposition after it performs the actual sector solve.

\begin{lemma}[Maximal rank and adic substitution]\label{arr:formal-replacement}
Let \(f_0=\mathcal F-\mathcal F(0)\) and
\(f_i=\partial_{p_i}\mathcal F\).
There are finitely many expressions \(g_1,\dots,g_r\), each
obtained from these formal series and coordinates by finitely
many partial differentiations and polynomial operations,
such that the following holds.
After choosing coordinates \(p=(u,v)\), \(g_u\) has a minor
\(\Delta\) nonzero along \(\widehat p\).
Replacing the free coordinates \(v\) by sufficiently accurate
polynomial approximations and solving \(g=0\) formally preserves
all the equations \(f_i=0\).  The substitution converges
adically despite possible divergence of \(\widehat p\).
\end{lemma}

\begin{proof}
Let \(\mathcal A\subset\mathbb R[[p]]\) be the full differential
polynomial algebra generated by the \(f_i\) and all coordinate
functions.  Let
\[
 I=\{f\in\mathcal A:f(\widehat p(s))=0\},\qquad
 K_s=\mathbb R((s)).
\]
Choose \(g_1,\ldots,g_r\in I\) with maximal gradient rank over
\(K_s\).  Although the closure \(\mathcal A\) is infinite,
each selected expression is finite and \(r\le\dim p\).
Split \(p=(u,v)\) so \(A=g_u\) is invertible along the arc,
and put \(\Delta=\det A\) and
\(m=\operatorname{ord}_s\Delta(\widehat p)<\infty\).
For each free coordinate set
\[
 D_j=\Delta\partial_{v_j}
       -\bigl(\operatorname{adj}(A)g_{v_j}\bigr)\cdot\partial_u,
 \qquad E_j=\Delta^{-1}D_j.
\]
The operators \(D_j\) preserve \(\mathcal A\), and
\(D_jg_i=0\) identically.
For every \(f\in I\), maximality implies
\(df(\widehat p)\in\operatorname{span}_{K_s}\{dg_i(\widehat p)\}\).
Hence \(D_jf(\widehat p)=0\), or \(D_jI\subset I\).
Ordinary partial differentiation need not preserve \(I\);
it is this cofactor tangency that does.

The denominators of iterated intrinsic derivatives admit an
explicit bound.  For \(n\ge1\) and \(h\in\mathcal A\),
\begin{equation}\label{arr:denominator}
 E_{j_1}\cdots E_{j_n}h
     =\frac{P_n}{\Delta^{2n-1}},\qquad P_n\in\mathcal A.
\end{equation}
The first denominator is \(\Delta\), and the induction is
\[
 E_j(P/\Delta^q)
   =\frac{\Delta D_jP-qP D_j\Delta}{\Delta^{q+2}}.
\]
If \(h\in I\), all numerators remain in \(I\).
Thus every iterated intrinsic derivative of every original
constraint vanishes along the arc.

Translate to \(\widehat p(s)\) and apply formal implicit
inversion over \(K_s\) to obtain a graph in
\(\eta=v-\widehat v(s)\).
On that graph \(E_j\) is differentiation in \(\eta_j\).
Apply \eqref{arr:denominator} also to the coordinate functions
\(u_i\).  Each coefficient of total \(\eta\)-degree \(n\)
has \(s\)-valuation at least \(-(2n-1)m\).
If \(\operatorname{ord}_s\eta\ge N>2m\), its degree \(n\)
contribution therefore has valuation at least
\[
 nN-(2n-1)m=m+n(N-2m)\longrightarrow\infty.
\]
The substitution is adically convergent and
the dependent displacement has order at least \(N-m\).
All original constraints have zero intrinsic Taylor
coefficients, so remain zero after substitution.
Both free and dependent displacements have positive
valuation; the same bound justifies rearranging their
formal evaluations.  This is not an evaluation of a
divergent numerical series.
If \(r=0\), take \(\Delta=1\) and all coordinates free;
the same argument says all relevant Taylor coefficients
vanish.  If there are no free coordinates, no substitution
is necessary.
\end{proof}

\begin{proposition}[A regular implicit solve in each array]
\label{arr:regular-implicit}
For the \(g_i\) selected above, form actual expressions
\(g_{\pm,i}(s,p)\) by applying the same finite differentiations
and polynomial operations to
\[
 G_\pm(s,p)-\mathcal F(0),\qquad \nabla_pG_\pm(s,p).
\]
There are parameter functions \(p_\pm(s)\) on the two
sectors, obtained by regular analytic implicit operations
after fixed-power rescaling, for which \(g_\pm(s,p_\pm)=0\).
They have the same formal expansion, with real coefficients.
For every original constraint
\[
 e_{\pm,0}=G_\pm(s,p_\pm)-\mathcal F(0),\qquad
 e_{\pm,i}=\partial_{p_i}G_\pm(s,p_\pm),
\]
there is \(c_2(D)>0\) such that
\begin{equation}\label{arr:initial-flatness}
 |e_{\pm,i}|\le C_D e^{-c_2(D)/|a|},\qquad
 |\bar\partial_s e_{\pm,i}|
       \le C_D|s|^{-M}e^{-D/(6|a|)}.
\end{equation}
\end{proposition}

\begin{proof}
First fix the inner evaluation ball of \eqref{arr:parameter-reserve},
with exponent \(L_0\), and the finite derivative orders occurring in
the selected \(g_i\).  Its polynomial distance from the outer solve
ball fixes all Cauchy bounds and possible polynomial losses before any
further rescaling.
Let \(m\) be the valuation of the selected minor.
If the second derivatives needed below are bounded, choose
\[
 \ell\ge\max\{L_0+1,m+1\}.
\]
If they have a previously fixed loss \(|s|^{-M_0}\), choose
instead \(\ell>\max\{L_0,m+M_0\}\).
In either case enlarge $\ell$ to the fixed threshold required in
Lemma~\ref{arr:gevrey}, in particular $\ell\ge k$, and then set
$N=\ell+m+1$.
Increase \(N\) and the polynomial approximation order if
needed to preserve all finite derivatives and leading
minor coefficients.
Take a real polynomial center \(p_0=(u_0,v_0)\) agreeing
with \(\widehat p\) modulo \(s^N\), and freeze \(v=v_0\).
In either array put \(J_0(s)=g_u(s,p_0(s))\).
The finite-order expansion and the nonzero formal minor give
\[
 g(s,p_0)=O(s^N),\quad
 \det J_0=s^m(d_0+O(s)),\quad d_0\ne0,\quad
 \|J_0^{-1}\|=O(|s|^{-m}).
\]
Allowing a larger fixed order for the center makes these
statements unchanged when a fixed derivative loss is present.

Set \(u=u_0+s^\ell h\), and consider the actual equation
\begin{equation}\label{arr:normalized-implicit}
 H(s,h)=s^{-\ell}J_0(s)^{-1}
       g(s,u_0+s^\ell h,v_0)=0 .
\end{equation}
Taylor's formula in \(u\) is the exact identity
\[
 H=h+s^{-\ell}J_0^{-1}g(s,p_0)
   +s^\ell J_0^{-1}\int_0^1(1-t)
     g_{uu}(s,u_0+ts^\ell h,v_0)[h,h]\,dt.
\]
Thus \(H(s,0)=O(s)\) and
\[
 H_h=I+O(|s|^{\ell-m-M_0})
\]
on a fixed \(h\)-ball, with \(M_0=0\) in the bounded case.
It is a uniformly regular implicit problem, and contraction
gives \(h=O(s)\).
Since $\ell>L_0$, the whole fixed $h$-ball used in this equation lies
strictly inside the previously chosen inner evaluation ball for small
$s$.  The physical displacement of its root is the still smaller
\(O(s^{\ell+1})\).
Free coordinates have been approximated more accurately
than the allowed dependent displacement; this pays the
amplification by \(s^{-m}\).

We check the divisions needed for Gevrey closure uniformly
in \(h\).  Write
\[
 H=\bigl(s^{-m}\det J_0\bigr)^{-1}
        s^{-(\ell+m)}
        \operatorname{adj}(J_0)
        g(s,u_0+s^\ell h,v_0).
\]
The first factor is the inverse of a unit.
The numerator of the second quotient has zero formal
coefficients below order \(\ell+m\), identically in \(h\):
the constant term has order \(N>\ell+m\), the linear
term is \(\det J_0\,s^\ell h\), and the quadratic term
starts at order \(2\ell\ge\ell+m\).
Hence Lemma~\ref{arr:division}, with its low-Taylor
subtractions for the disk approximants, applies to these
actual quotients.  Their leading equation is \(H_0(h)=h\).
The regular implicit roots have the Gevrey expansion
given by formal inversion.
Lemma~\ref{arr:formal-replacement} identifies it with the
substituted formal arc and proves that every original
constraint has zero formal expansion.  This gives the
first estimate in \eqref{arr:initial-flatness}.

For the antiholomorphic estimate differentiate the exact
selected equations.  Holomorphy in \(p\) and the polynomial
nature of \(v_0\) give
\begin{equation}\label{arr:antiholomorphic-chain}
 \bar\partial_s u
    =-g_u(s,p(s))^{-1}\bar\partial_sg(s,p(s)),\qquad
 \bar\partial_s f(s,p(s))
    =\bar\partial_sf+f_u\,\bar\partial_su .
\end{equation}
Here \(H_h=J_0^{-1}g_u\) is uniformly invertible, so the
actual \(g_u^{-1}\) has only the fixed loss
\(O(|s|^{-m})\).
Every selected \(g_i\) uses finitely many parameter
derivatives and polynomial operations, so
\eqref{arr:dbar} bounds its fixed-parameter
antiholomorphic derivative with only further fixed powers.
The chain rule therefore loses only powers of \(|s|^{-1}\).
For fixed \(D\) these are absorbed between the constants
\(D/5\) and \(D/6\), proving the second bound.
The constants and the ball radius \(c_D|s|^{L_0}\) may
depend on \(D\); no exponentially small ball is used.
\end{proof}

\subsection{Improving flatness on the evaluation ray}

At this point each action's difference from the base critical value,
and its gradient, are exponentially small at their separately selected
parameters, but the decay constant is not yet
large enough for the jump comparison.  Their much smaller
antiholomorphic derivatives allow the sector argument below to improve
that constant.  We then move both actions to one common parameter value.

\begin{lemma}[Almost-holomorphic exponential improvement]
\label{arr:flatness-upgrade}
Suppose an error \(e(w)\) on an exterior sector
\[
 |w|>R,\qquad |\arg w|<\pi/2+\delta
\]
satisfies
\[
 |e(w)|\le C_D e^{-c_2(D)|w|},\qquad
 |\bar\partial_we(w)|\le C_D(1+|w|)^M e^{-D|w|/6},
 \quad c_2(D)>0.
\]
Then on the ray \(\arg w=\pi/2\),
\begin{equation}\label{arr:improved-exponent}
 |e(w)|\le C'_D
       \exp\{-D\sin(\delta/2)|w|/8\}.
\end{equation}
In particular, if \(\delta\ge cD^{-1/2}\), the obtainable
decay coefficient tends to infinity with \(D\).
\end{lemma}

\begin{proof}
Put \(A=(D/8)e^{-i\delta/2}\).
Choose \(\epsilon>0\) so small that
\[
 \epsilon<\delta/2,\qquad
 D\sin\epsilon/8<c_2(D)/2.
\]
For each fixed \(D\) both conditions can be met, regardless
of how small \(c_2(D)\) is.
The sector between the rays
\[
 \alpha=\delta/2-\pi/2+\epsilon,\qquad
 \beta=\delta/2+\pi/2-\epsilon
\]
has opening \(\pi-2\epsilon<\pi\), lies within the
available sector, and contains the evaluation ray strictly.
On its two sides \(\operatorname{Re}(Aw)
=(D/8)|w|\sin\epsilon\), so \(H=e^{Aw}e\) is bounded.
It also has at most order-one exponential growth and is
bounded on the finite inner circular boundary.
Its antiholomorphic source satisfies
\[
 |\bar\partial H|\le C_D(1+|w|)^M e^{-D|w|/24}.
\]

Extend this source, not \(H\), by zero to the plane, calling
the extension \(Q\).  It belongs to \(L^1\cap L^\infty\).
The Cauchy transform
\[
 V(w)=\frac1\pi\int_{\mathbb C}
             \frac{Q(\zeta)}{w-\zeta}\,dA(\zeta)
\]
is bounded.  Indeed the contribution from
\(|w-\zeta|<1\) is at most \(2\|Q\|_\infty\), and the
complement is at most \(\pi^{-1}\|Q\|_1\).
Distributionally \(\bar\partial V=Q\); hence \(H-V\) is
holomorphic in the exterior sector, has bounded boundary
values including the inner arc, and has order at most one.

For completeness the required Phragm\'en--Lindel\"of step
follows from the maximum principle.  Choose
\[
 1<\rho<\frac{\pi}{\pi-2\epsilon}.
\]
With sector center \(\theta_0=\delta/2\), the real part of
\((e^{-i\theta_0}w)^\rho\) is positive on the closed sector.
Multiply \(H-V\) by
\(\exp\{-\tau(e^{-i\theta_0}w)^\rho\}\), \(\tau>0\).
Its outer-circle values tend to zero as the radius tends
to infinity, since \(\rho>1\); its side and inner-arc
values are bounded by the original boundary bound.
The maximum principle on truncated sectors followed by
the limits at infinity and then \(\tau\downarrow0\)
shows \(H-V\) bounded.  Since \(V\) is bounded, so is \(H\).
At \(\arg w=\pi/2\),
\(\operatorname{Re}(Aw)=D|w|\sin(\delta/2)/8\),
which proves \eqref{arr:improved-exponent}.
\end{proof}

For the upper array use \(w=i/a=i/s^k\).
Its retained sector has the geometry of the lemma with excess
$\delta_{\rm use}\ge cD^{-1/2}$; changing from
\(s\) to this branch of \(w\) adds only a polynomial
factor to the antiholomorphic derivative.
Apply Lemma~\ref{arr:flatness-upgrade} to each error in
\eqref{arr:initial-flatness}.  The lower array is treated
by reflection of the sector variable.  Thus, for any
prescribed finite \(K>0\), first choosing \(D\) sufficiently
large and then \(s\) sufficiently small gives
\begin{equation}\label{arr:arbitrary-flatness}
 |G_\pm(s,p_\pm)-\mathcal F(0)|
       +|\nabla_pG_\pm(s,p_\pm)|
      \le C_K e^{-K/|a|}
\end{equation}
on the evaluation ray.  The \(D\)-dependent constants
do not affect this order of limits.

\begin{proposition}[The common-parameter action upper bound]
\label{arr:action-upper}
The formal critical arc \eqref{arr:critical-arc} implies
that there are common parameters \(p(s)\), having a formal
expansion with real coefficients, such that
\begin{equation}\label{arr:upper-estimate}
 G_+(s,p(s))-G_-(s,p(s))
                  =O(e^{-1.25A_{\mathrm{lead}}}).
\end{equation}
Here \(A_{\mathrm{lead}}\) is the minimum leading adverse action scale
in \eqref{arr:A0}.  If there is a negative real leading
opening, the evaluation is on positive real \(s\).
If all nonzero leading openings are positive, the
evaluation is at \(a<0\) with the two continuations taken
on the same sheet in \(s\); the common parameters may
then be complex.  In particular this conclusion includes
even \(k\).  After a sufficiently small fixed generic
change of evaluation angle the same assertion holds with
\(1.2\) in place of \(1.25\), using
\(A_{\mathrm{lead}}=\min |Z_{j,\mathrm{lead}}|^2/4\).
\end{proposition}

\begin{proof}
Use the same polynomial center and the same free
coordinates for the two regular implicit solves.
The finite expressions defining \(g_\pm\), the selected
inverse minor and all rescalings have only polynomial
losses.  Apply \eqref{arr:fixed-parameter-comparison} on
their common inner evaluation ball and then subtract the
two uniformly contracting normalized equations.
It follows that
\begin{equation}\label{arr:parameter-comparison}
 |p_+(s)-p_-(s)|
      \le C_D|s|^{-M}e^{-(0.7-o(1))A_{\mathrm{lead}}}.
\end{equation}
Their midpoint \(p=(p_++p_-)/2\) lies in the same ball, with room
provided by the fixed-$h$ inclusion proved above.
When the data and evaluation parameter are real, choose
conjugate contours and prescriptions; uniqueness gives
\(p_-=\overline{p_+}\), so the midpoint is real.
The estimate itself uses no conjugacy.

The Hessians of the actions have polynomial bounds on
the common ball.  Taylor expansion from each \(p_\pm\)
to \(p\), with \eqref{arr:arbitrary-flatness}, gives
\[
 |G_\pm(s,p)-\mathcal F(0)|
 \le C_K e^{-K/|a|}
      +C_K e^{-K/|a|}|p-p_\pm|
      +C_D|s|^{-M}|p-p_\pm|^2.
\]
Since \(A_{\mathrm{lead}}\le C_A/|a|\) for a fixed $C_A$, choose
$K>1.4C_A$, so $K/|a|>1.4A_{\mathrm{lead}}$.
The last term has exponent
\((1.4-o(1))A_{\mathrm{lead}}\); the fixed powers are absorbed in the
margin between \(1.4\) and \(1.25\).
Subtracting the two estimates proves
\eqref{arr:upper-estimate}.
The common formal parameter series has real coefficients
because the selected formal equations and the polynomial
free path are real and formal implicit inversion is unique.

If all nonzero \(c_j\) are positive, changing the sign
of real \(s\) does not solve the problem when \(k\) is
even.  Instead choose the continuous argument with
\(\arg a=\pi\), so \(\arg s=\pi/k\).
Use the two argument intervals
\((-\delta_{\rm use},\pi+\delta_{\rm use})\) and
\((\pi-\delta_{\rm use},2\pi+\delta_{\rm use})\) for \(a\), taking
\(\arg s=(\arg a)/k\) continuously on both.  Their
overlap therefore gives the same value of \(s\) near
\(\arg s=\pi/k\), not values differing by \(2\pi/k\).
Equivalently, all the preceding constructions are made
on a fixed local covering chart in the \(s\)-plane;
\(a\) only labels its leading scale.
The effective opening coefficients are now negative
real to leading order.  Higher terms need not be real.
Strict contour margins persist under these \(o(1)\)
changes, and the common-curtailment estimate has only
\(1+o(1)\) changes in its Gaussian heights.
The flatness-upgrade variable is rotated accordingly.
Thus \eqref{arr:parameter-comparison} and the same
midpoint Taylor estimate hold with complex \(p\).
No identification of the two arrays by conjugation,
and no reality of their actual subleading parameters,
has been assumed.

Finally fix \(D\), all contour margins, and a decay
constant in \eqref{arr:arbitrary-flatness} with room to
spare.  A sufficiently small fixed change of evaluation
angle stays inside the sectors and preserves the
positive multiplier gain in
Lemma~\ref{arr:flatness-upgrade}.
The endpoint Gaussian real parts change continuously
by a small relative amount.  The preceding \(1.4\)
quadratic margin therefore still gives exponent \(1.2\).
This is the open range of angles available for the
later separation of distinct Laurent principal parts.
\end{proof}
